\chapter{Demonstration Runbook: Install to Uninstall}
\label{ch:runbook}
This is the whole demonstration, in order, from an empty host to an empty host. Follow it top to bottom.
Every step has the command, what you should see, and an honest status tag:
\begin{center}\small
\begin{tabular}{@{}ll@{}}
\verified & the step was run on the Ubuntu 24.04 test VM and behaved as written\\[3pt]
\notrun & the step has not been run; the expected output is derived from the code\\[3pt]
\fromsrc & as above, and the exact wording is taken from the source rather than a live run\\
\end{tabular}
\end{center}
Steps marked \notrun\ are the ones to rehearse before an audience. Section~\ref{sec:rb-auto} shows how to run
the install and checks automatically on the VM.

\section{Time budget and layout}
\begin{center}\small
\begin{tabularx}{\textwidth}{@{}L{1.2cm}L{4.4cm}Y@{}}
\toprule
\textbf{Phase} & \textbf{What} & \textbf{Screen / terminal}\\
\midrule
0--3 & Prepare, install, verify (about 4 min) & Terminal A (the demo host)\\
4 & Start the stack (1 min) & Terminal A\\
5 & Sensor and control plane live (2 min) & Terminal B: \texttt{journalctl} follow\\
6 & Web dashboard (2 min) & Browser\\
7 & SSH console (2 min) & Terminal C\\
8 & Contain a process (2 min) & Terminals A and C\\
9--10 & Audit proof, policy reload (2 min) & Terminal A\\
11 & Failure and recovery (1 min) & Terminals A and C\\
12--14 & Stop, uninstall, reset (3 min) & Terminal A\\
\bottomrule
\end{tabularx}
\end{center}
Open three terminals on the demo host before you start (A: commands, B: logs, C: the SSH console) and one browser window.

\section{Phase 0: Before you start}
\resetsteps
\step{Confirm the host meets the prerequisites (Chapter~\ref{ch:prereq})}
\begin{lstlisting}[style=console]
$ ls /sys/kernel/btf/vmlinux            # must exist
$ cat /sys/kernel/security/lsm          # must end with: ...,bpf
$ id -nG | tr ' ' '\n' | grep -x -E 'sudo|wheel'   # you need sudo
\end{lstlisting}
\verified\ (\texttt{lockdown,capability,landlock,yama,apparmor,bpf} and BTF present on the test VM.)
\step{Make sure no earlier install is running}
\begin{lstlisting}[style=console]
$ systemctl is-active kbd kb-sensor kbopd sshd@kb-operator kb-checker
\end{lstlisting}
All should say \texttt{inactive} (or \texttt{unknown} on a clean host). \verified\ If a development install is running, either
stop it by hand or let the package's migration step retire it (Section~\ref{sec:migrate}).
\step{Check the package file}
\ifarch
\begin{lstlisting}[style=console]
$ sha256sum -c SHA256SUMS --ignore-missing
kernel-borderlands-1.1.0-1-x86_64.pkg.tar.zst: OK
\end{lstlisting}
\else
\begin{lstlisting}[style=console]
$ sha256sum -c SHA256SUMS --ignore-missing
kernel-borderlands_1.1.0_amd64.deb: OK
\end{lstlisting}
\fi
\verified

\section{Phase 1: Install}
\resetsteps
\step{Install the package}
\ifarch
\begin{lstlisting}[style=console]
$ sudo pacman -U kernel-borderlands-1.1.0-1-x86_64.pkg.tar.zst
\end{lstlisting}
\notrun\ (not run on Arch; the post-install hook is the same script as on Ubuntu.) Expected tail:
\begin{lstlisting}[style=console]
created group: kb
created user: operator (member of kb)
generated SSH host key: /etc/kb/ssh/ssh_host_ed25519_key
created empty /etc/kb/authorized_keys - add operator public keys here
provisioning complete
kb: generated the dashboard API token in /etc/kb/api.env
==> Kernel Borderlands installed. Nothing has been started.
\end{lstlisting}
\else
\begin{lstlisting}[style=console]
$ sudo apt install ./kernel-borderlands_1.1.0_amd64.deb
\end{lstlisting}
\verified\ Expected tail (on a clean host):
\begin{lstlisting}[style=console]
Setting up kernel-borderlands (1.1.0) ...
created group: kb
created user: operator (member of kb)
generated SSH host key: /etc/kb/ssh/ssh_host_ed25519_key
created empty /etc/kb/authorized_keys - add operator public keys here
provisioning complete
kb: generated the dashboard API token in /etc/kb/api.env
\end{lstlisting}
On a host that had the old development install you will also see \texttt{kb: moved legacy ...}, \texttt{kb: legacy dev unit ... moved aside},
\texttt{kb: operator shell ... -> /bin/sh} and \texttt{kb: operator account was locked ...}. These are the migration working. \verified
\fi
\step{Prove nothing started and the checker is held down}
\begin{lstlisting}[style=console]
$ kbctl status
kbd.service                  inactive
kb-sensor.service            inactive
kbopd.service                inactive
sshd@kb-operator.service     inactive
kb-checker.service           UNDER MAINTENANCE AND REPAIRS (kept down)
$ systemctl is-enabled kb-checker
disabled
\end{lstlisting}
\verified
\step{Check the provisioning}
\begin{lstlisting}[style=console]
$ getent passwd operator
operator:x:998:994::/:/bin/sh
$ sudo stat -c '%a %U:%G' /etc/kb/api.env
640 root:kb
\end{lstlisting}
\verified\ (UID and GID numbers differ per host; the home is \texttt{/} and the shell \texttt{/bin/sh}.)

\section{Phase 2: Authorize an operator and start the stack}
\resetsteps
\step{Authorize your SSH public key for the console}
\begin{lstlisting}[style=console]
$ sudo kbctl ssh add-key ~/.ssh/id_ed25519.pub
key authorized for operator
\end{lstlisting}
\verified\ (If you have no key: \texttt{ssh-keygen -t ed25519}.)
\step{Bring the stack up}
\begin{lstlisting}[style=console]
$ sudo kbctl up
$ kbctl status
kbd.service                  active
kb-sensor.service            active
kbopd.service                active
sshd@kb-operator.service     active
kb-checker.service           UNDER MAINTENANCE AND REPAIRS (kept down)
\end{lstlisting}
\verified
\step{Check the sockets and ports}
\begin{lstlisting}[style=console]
$ ls -l /run/kb
srw-rw---- 1 root kb   ... kba.sock       # group kb: the operator console can reach it
srwxr-xr-x 1 root root ... kbct.sock
srwxr-xr-x 1 root root ... kbd.sock
$ ss -ltn | grep -E ':(8080|8082|2222)\b'
\end{lstlisting}
\verified\ Three listeners: \texttt{127.0.0.1:8080} (\cmd{kbd} API), \texttt{127.0.0.1:8082} (dashboard), \texttt{0.0.0.0:2222} (operator \cmd{sshd}).

\section{Phase 3: Show the sensor and the control plane working}
\resetsteps
\step{In Terminal B, follow the sensor and \cmd{kbd}}
\begin{lstlisting}[style=console]
$ sudo journalctl -f -u kb-sensor -u kbd
\end{lstlisting}
Since 1.1.0 the sensor is \emph{quiet by default}: it logs startup, the occasional state transition, and one summary line every 30 seconds
(events seen, lines suppressed, its own syscalls skipped in the kernel, critical self-activity). It no longer prints one line per event: on a busy
host that flooded the journal. To see the old per-event stream while demonstrating, turn it on (it is capped at 200 lines per second):
\begin{lstlisting}[style=console]
$ sudo systemctl edit kb-sensor      # add:  [Service]  Environment=KB_SENSOR_VERBOSE=1
$ sudo systemctl restart kb-sensor
\end{lstlisting}
Remove the override afterwards. With it on you will see \texttt{[process\_exec]}, \texttt{[process\_exit]} and \texttt{[privilege\_change]} lines.
\verified
\step{In Terminal A, show the live counters}
\begin{lstlisting}[style=console]
$ kbctl stats
events/sec:       112.40
active processes: 238
\end{lstlisting}
\fromsrc\ (numbers will differ).
\step{Trigger a detection with a safe, real event}
Reading a credential file is something the sensor always flags. Do it from Terminal A:
\begin{lstlisting}[style=console]
$ sudo cat /etc/shadow > /dev/null
\end{lstlisting}
In Terminal B you should see lines like these (observed on the VM with the 1.1.0 defaults, produced by \cmd{sudo} and \cmd{cat} themselves):
\begin{lstlisting}[style=console]
[BEHAVIOR ENGINE] PID=60970 COMM=sudo State transition: OBSERVED -> SUSPICIOUS (Reason: CRED_FILE_ACCESS, ...)
[BEHAVIOR ENGINE] PID=60972 COMM=cat  State transition: OBSERVED -> SUSPICIOUS (Reason: CRED_FILE_ACCESS, ...)
[KB] Zone PID=60970 COMM= SAFE->SUSPICIOUS score=18.0
[AUDIT] ZONE_SAFE_TO_SUSPICIOUS | pid=60970 comm= score=18.03 ...
\end{lstlisting}
The raw \texttt{[file\_access] ... file=/etc/shadow} line appears only with \texttt{KB\_SENSOR\_VERBOSE=1}.
\verified\ Say what happened: the kernel hook saw the file open, the behavior engine moved the process to a higher state, and
\cmd{kbd} recorded the zone change in the tamper-evident audit log. This is a \emph{detection}, not a block: containment is an operator
decision (Phase 6).
\begin{kbnote}[WHY THERE IS NO ATTACK SCRIPT HERE]
The repository's \file{scripts/attack-lab/} README lists six attack simulations, but only \file{containment_demo.py} exists in the
tree. This runbook therefore uses only harmless, real events: a credential-file read and a \texttt{sleep} process.
Do not run attack simulations on a host that is not isolated.
\end{kbnote}

\section{Phase 4: Web dashboard}
\resetsteps
\step{Open it with the API token handed over}
\begin{lstlisting}[style=console]
$ kbctl dashboard open
opened http://127.0.0.1:8082
\end{lstlisting}
\verified\ on the VM: the dashboard service is active and \texttt{/} returns 200. \notrun\ for \cmd{kbctl dashboard open} itself on the VM (it was tested locally against the packaged \cmd{kbopd} with a stub \cmd{systemctl}) and for the page content in a browser. The page loads at \texttt{http://127.0.0.1:8082},
the token is moved from the URL fragment into \texttt{sessionStorage}, and the address bar is cleaned. On a headless server, run
\texttt{kbctl dashboard url --token} on the host and browse through an SSH tunnel (Chapter~\ref{ch:dash}).
\step{Show the data and the action buttons}
Point at the process table, the zone chart and the alert feed. Then show that the action buttons are authorized:
\begin{lstlisting}[style=console]
$ TOK=$(sudo grep ^KB_HTTP_API_TOKEN= /etc/kb/api.env | cut -d= -f2)
$ curl -s -o /dev/null -w '%{http_code}\n' -X POST http://127.0.0.1:8080/api/policy/reload                         # no token
401
$ curl -s -o /dev/null -w '%{http_code}\n' -X POST -H "Authorization: Bearer $TOK" http://127.0.0.1:8080/api/policy/reload
200
\end{lstlisting}
\verified\ (the 401 / 401 (wrong token) / 200 behaviour of the API). Clicking isolate and restore in the page is \notrun.

\section{Phase 5: The SSH console}
\resetsteps
\step{In Terminal C, connect}
\begin{lstlisting}[style=console]
$ kbctl ssh connect
\end{lstlisting}
\verified\ Key accepted, the terminal UI appears with four tabs (\emph{Processes}, \emph{Alerts}, \emph{Agent Activity}, \emph{Query Console}), live link to \cmd{kbd}
(no \texttt{[OFFLINE]} banner), and it exits cleanly on \texttt{q}. From another machine: \texttt{ssh -t -p 2222 operator@HOST}.
\step{Drive it}
\texttt{Tab} through the views; on \emph{Processes} press \texttt{/} and type \texttt{sleep}; on \emph{Query Console} try
\texttt{help}, \texttt{stats}, \texttt{zone suspicious}. The tab rendering is \verified; the individual query commands are \notrun.
\step{Show that it is a forced console, not a shell}
\begin{lstlisting}[style=console]
$ ssh -p 2222 operator@HOST id
\end{lstlisting}
You get the console, not the output of \texttt{id}. \notrun

\section{Phase 6: Contain a process}
\resetsteps
\step{Start a harmless victim in Terminal A}
\begin{lstlisting}[style=console]
$ sleep 600 &
[1] 5120
\end{lstlisting}
\step{Contain it from the command line}
\begin{lstlisting}[style=console]
$ sudo kbctl process isolate --pid 5120 --level CGROUP --reason "demo"
pid=5120 contained at level CGROUP
\end{lstlisting}
\fromsrc\ The output format is taken from the source. Escalation order: \texttt{CGROUP}, \texttt{SECCOMP}, \texttt{NAMESPACE}, \texttt{TERMINATE}.
\step{Show it in the console}
In Terminal C: \emph{Processes}, filter \texttt{sleep}; the \textsc{Containment} column shows the level. Or contain from the TUI: select the row,
press \texttt{c}, choose a level (\texttt{1}--\texttt{4}), confirm with \texttt{y}. \notrun
\step{Release and clean up}
\begin{lstlisting}[style=console]
$ kill 5120
\end{lstlisting}

\section{Phase 7: The audit proof}
\resetsteps
\step{Verify the hash chain}
\begin{lstlisting}[style=console]
$ kbctl audit verify
chain intact - 1763 entries verified
\end{lstlisting}
\verified\ (the count differs per host). Every record hashes the one before it; one altered byte breaks the chain.
\step{Export a copy}
\begin{lstlisting}[style=console]
$ kbctl audit export --out /tmp/demo-audit.json
wrote 1763 entries to /tmp/demo-audit.json
\end{lstlisting}
\fromsrc

\section{Phase 8: Change policy live}
\resetsteps
\step{Edit and reload}
\begin{lstlisting}[style=console]
$ sudo $EDITOR /etc/kb/policy.yaml      # e.g. lower bash's suspicious threshold
$ sudo kbctl policy reload
\end{lstlisting}
\notrun\ for \cmd{kbctl policy reload} (gRPC). The same reload through the HTTP API with the token is \verified\ (Phase~4).

\section{Phase 9: Failure and recovery}
\resetsteps
\step{Stop the control plane under the console}
\begin{lstlisting}[style=console]
$ sudo systemctl stop kbd
\end{lstlisting}
In Terminal C the TUI shows the \texttt{[OFFLINE] kbd unreachable} banner and fictional \texttt{[DEMO DATA]} rows. Point out that the banner is the
honest signal and the rows are not live. \notrun\ for stopping \cmd{kbd} under a live console. The banner itself was seen on the VM when the link to \cmd{kbd} failed (before the socket-group fix) and in unprivileged runs.
\step{Bring it back}
\begin{lstlisting}[style=console]
$ sudo kbctl up
$ kbctl status
\end{lstlisting}
\cmd{kbctl up} after a stop is \verified. Whether an already-open console reconnects on its own is \notrun; if it does not, quit with \texttt{q} and run \texttt{kbctl ssh connect} again.

\section{Phase 10: Optional, the agent swarm}
Only if the demo audience wants it, and only on a host where you have built the virtualenv (Chapter~\ref{ch:aads}).
\begin{lstlisting}[style=console]
$ sudo systemctl start kbagents
$ curl -s http://127.0.0.1:8601/agents | head -c 300
\end{lstlisting}
Running the launcher by hand against an existing venv is \verified\ (nine agents, zero errors); running it as the systemd unit is \notrun.

\section{Phase 11: Stop everything}
\begin{lstlisting}[style=console]
$ sudo kbctl down
$ kbctl status
kbd.service                  inactive
...
kb-checker.service           UNDER MAINTENANCE AND REPAIRS (kept down)
\end{lstlisting}
\verified\ (\cmd{kbctl down} stops in reverse boot order.) To also stop them from starting at boot:
\texttt{sudo kbctl dashboard down; sudo kbctl ssh down}.

\section{Phase 12: Uninstall}
\resetsteps
\step{Remove the package}
\ifarch
\begin{lstlisting}[style=console]
$ sudo pacman -R kernel-borderlands
\end{lstlisting}
\else
\begin{lstlisting}[style=console]
$ sudo apt remove kernel-borderlands        # keeps config you changed
$ sudo apt purge  kernel-borderlands        # also removes config files the package owns
\end{lstlisting}
\fi
\notrun\ for the removal itself (the pre-removal hook disables and stops the units). The service-stopping part is the same logic as
\texttt{kbuninstaller.sh}.
\step{See what is deliberately left behind}
\begin{lstlisting}[style=console]
$ ls /etc/kb /var/lib/kbd
$ getent passwd operator; getent group kb
$ ls /etc/systemd/system/*.pre-migration 2>/dev/null      # only if a dev install was migrated
\end{lstlisting}
These hold keys, the token, the audit chain and the account; a package manager should not destroy them.
\step{Wipe everything (only if you want a truly clean host)}
\begin{lstlisting}[style=console]
$ sudo rm -rf /etc/kb /var/lib/kbd /opt/kb
$ sudo rm -f /etc/systemd/system/*.pre-migration
$ sudo userdel operator && sudo groupdel kb
$ sudo systemctl daemon-reload
\end{lstlisting}
\notrun
\step{Confirm}
\begin{lstlisting}[style=console]
$ ss -ltn | grep -E ':(8080|8082|2222)\b' || echo "ports free"
$ command -v kbd kbctl || echo "binaries gone"
\end{lstlisting}
\notrun

\section{Phase 13: Reset between demonstrations}
To repeat the demonstration from nothing, either do Phase~12 and Phase~1 again, or use the two scripts (next section) which do the same on the Ubuntu VM.

\section{Running install and checks automatically}
\label{sec:rb-auto}
The repository ships two scripts in \file{packaging/vm-test/}; they prompt for \cmd{sudo} once and must be run with a terminal.
\begin{lstlisting}[style=console]
$ ssh -t HOST 'bash ~/kbuninstaller.sh --remove-package'   # stop services, remove the package (keeps data)
$ ssh -t HOST 'bash ~/kbinstaller.sh'                       # install, restart, check everything
...
=== SUMMARY ===
PASS=34 FAIL=0
\end{lstlisting}
\verified\ \file{kbinstaller.sh} checks Phases 1--5 and 7 (install, migration, token, CORS and auth, dashboard, services, the SSH console with a live
\cmd{kbd} link, audit verify, sensor). It does \emph{not} cover contain, policy reload, failure and recovery, the swarm, reboot persistence or
uninstall; those are the \notrun\ rows above.

\section{The whole demonstration on one page}
\ifarch
\begin{lstlisting}[style=console]
# --- install ---
sudo pacman -U kernel-borderlands-1.1.0-1-x86_64.pkg.tar.zst
\end{lstlisting}
\else
\begin{lstlisting}[style=console]
# --- install ---
sudo apt install ./kernel-borderlands_1.1.0_amd64.deb
\end{lstlisting}
\fi
\begin{lstlisting}[style=console]
kbctl status                                   # all inactive, checker labelled
# --- start ---
sudo kbctl ssh add-key ~/.ssh/id_ed25519.pub
sudo kbctl up && kbctl status
# --- show ---
sudo journalctl -f -u kb-sensor -u kbd         # terminal B
sudo cat /etc/shadow > /dev/null               # trigger a detection
kbctl dashboard open                           # browser
kbctl ssh connect                              # terminal C
sleep 600 & sudo kbctl process isolate --pid $! --level CGROUP --reason demo
kbctl audit verify
sudo kbctl policy reload
# --- recover ---
sudo systemctl stop kbd ; sudo kbctl up
# --- stop and remove ---
sudo kbctl down
\end{lstlisting}
\ifarch
\begin{lstlisting}[style=console]
sudo pacman -R kernel-borderlands
sudo rm -rf /etc/kb /var/lib/kbd ; sudo userdel operator ; sudo groupdel kb
\end{lstlisting}
\else
\begin{lstlisting}[style=console]
sudo apt purge kernel-borderlands
sudo rm -rf /etc/kb /var/lib/kbd ; sudo userdel operator ; sudo groupdel kb
\end{lstlisting}
\fi
